\documentclass[10pt]{article}
\usepackage[T1]{fontenc}
\usepackage[utf8]{inputenc}
\usepackage[french]{babel}
\usepackage{amsmath,amssymb,mathrsfs,mathpazo}
\usepackage{geometry}
\geometry{a4paper,margin=2cm}
\usepackage{graphicx,array,tabularx,tabularray,colortbl,multirow,multicol,booktabs,enumitem,pifont,ulem,fancybox,diagbox,cancel}
\usepackage{pstricks,pstricks-add,pst-plot,pst-tree,pst-node,pst-text,pst-func,pst-math,pst-bspline,pst-3dplot,pst-eucl,pst-solides3d}
\let\pscalculate\undefined
\usepackage{pst-bezier}
\usepackage{tikz,pgfplots,tkz-tab}
\usetikzlibrary{arrows,arrows.meta,patterns,calc,positioning,trees,decorations.markings,shapes.geometric,shapes.misc,shapes.arrows,intersections,decorations.pathmorphing,decorations.pathreplacing,backgrounds,shadows,babel}
\pgfplotsset{compat=1.18}
\usepgfplotslibrary{fillbetween}
\usepackage[np]{numprint}
\usepackage{listings,makecell,pas-tableur,twoopt}
\let\barre\undefined
\usepackage[pstricks]{tablvar}
\definecolor{darkgreen}{RGB}{0,100,0}
\providecommand{\red}{\color{red}}
\providecommand{\blue}{\color{blue}}
\providecommand{\green}{\color{green}}
\providecommand{\black}{\color{black}}
\providecommand{\white}{\color{white}}
\providecommand{\gray}{\color{gray}}
\providecommand{\cyan}{\color{cyan}}
\providecommand{\magenta}{\color{magenta}}
\providecommand{\yellow}{\color{yellow}}
\providecommand{\violet}{\color{violet}}
\providecommand{\orange}{\color{orange}}
\providecommand{\euro}{\texteuro}
\providecommand{\EUR}{\texteuro}
\providecommand{\ieme}{\textsuperscript{e}}
\providecommand{\degres}{\ensuremath{^\circ}}
\providecommand{\degree}{\ensuremath{^\circ}}
\providecommand{\vv}[1]{\overrightarrow{#1}}
\providecommand{\Coord}[1]{#1}
\providecommand{\decosix}{}
\providecommand{\starredbullet}{\ensuremath{\bullet}}
\pagestyle{empty}
\setlength{\parindent}{0pt}
\usepackage{adjustbox}\DeclareUnicodeCharacter{2236}{:}\DeclareUnicodeCharacter{2190}{\ensuremath{\leftarrow}}
\providecommand{\R}{\mathbb{R}}
\providecommand{\N}{\mathbb{N}}
\providecommand{\Z}{\mathbb{Z}}
\providecommand{\Q}{\mathbb{Q}}
\providecommand{\D}{\mathbb{D}}
\providecommand{\C}{\mathbb{C}}
\providecommand{\e}{{\mathrm{e}}}
\providecommand{\vect}[1]{\overrightarrow{#1}}
\providecommand{\vectt}[1]{\overrightarrow{\mathrm{#1}}}
\providecommand{\vvt}[1]{\overrightarrow{\mathrm{#1}}}
\providecommand{\np}[2][]{{#2~#1}}
\providecommand{\vv}[1]{{\overrightarrow{#1}}}
\renewcommand{\d}{{\mathrm{d}}}
\providecommand{\degres}{{^\circ}}
\providecommand{\degre}{{^\circ}}
\providecommand{\blue}{\color{blue}}
\providecommand{\red}{\color{red}}
\providecommand{\enskip}{\quad}
\providecommand{\llbracket}{\lbrack\!\lbrack}
\providecommand{\rrbracket}{\rbrack\!\rbrack}
\providecommand{\euro}{{\text{€}}}
\providecommand{\no}{{\text{n°}}}
\providecommand{\interval}[2]{{[#1;#2]}}
\providecommand{\intervFO}[2]{{[#1;#2[}}
\providecommand{\dotfill}{\ldots}

\providecommand{\ds}{\displaystyle}
\providecommand{\textding}[1]{#1}
\providecommand{\Oij}{\(\mathrm{O};\vec{\imath},\vec{\jmath}\)}
\providecommand{\Oijk}{\(\mathrm{O};\vec{\imath},\vec{\jmath},\vec{k}\)}
\providecommand{\Ouv}{\(\mathrm{O};\vec{u},\vec{v}\)}
\providecommand{\intervFF}[2]{\ensuremath{[#1;#2]}}
\renewcommand{\euro}{\ifmmode\text{\texteuro}\else\texteuro\fi}
\DeclareUnicodeCharacter{00B0}{\ifmmode^{\circ}\else\textdegree\fi}
\DeclareUnicodeCharacter{20AC}{\ifmmode\text{\texteuro}\else\texteuro\fi}
% Énoncés : Éducation nationale. Transcription et corrigés : APMEP, Denis Vergès et contributeurs. Découpage et classement : CoopMaths, Sébastien Lozano, Éric Elter et contributeurs. Adaptation de la présentation : Maths Explorer.
% CC BY-SA : https://coopmaths.fr/www/about/
% https://coopmaths.fr/alea/static/bac/2023/tex/bac_2023_09_sujet1_ameriquesud_2_cor.tex
\begin{document}
\noindent\textit{APMEP / CoopMaths. CC BY-SA. Présentation adaptée par Maths Explorer.}\par\medskip
\medskip

\begin{enumerate}
\item 

	\begin{enumerate}
		\item 
		On a $\dfrac{\np{4921}}{\np{18221965}} \approx \np{0,0161} \approx 1,6\,\%$.
		\item 
De même $\dfrac{\np{4921}}{\np{18221965}} \approx \np{0,0003}< 0,001$ soit moins de 0,1\,\%.













	\end{enumerate}
\item 




	\begin{enumerate}
		\item 
Les 20 accouchements sont des évènements indépendants et la probabilité d'avoir une naissance est égale à 0,016, donc la variable $X$ suit une loi binomiale 

$\mathcal{B}(20~;~0,016)$.

On a $p(X = 1) = 20 \times 0,016 \times 0984^{19} \approx \np{0,2355}$ soit 0,236 au millième près.
		\item 
		
On a $p(X \geqslant 1) = 1 - p(X = 0) = 1 - 0,984^n \times 0,016^0  = 1 - 0,984^n$.

Il faut donc résoudre l'inéquation :

$1 - 0,984^n \geqslant 0,99 \iff 0,01 \geqslant 0,984^n \iff$ (par croissance de la fonction logarithme népérien) $\ln 0,01 \geqslant n \ln 0,984 \iff n \geqslant \dfrac{\ln 0,01}{ \ln 0,984}$ soit $n \geqslant 285,5$. Donc en moyenne, sur 286 accouchements il y a un accouchement double.


	\end{enumerate}
\item Dans cette maternité, parmi les naissances double, on estime qu'il y a 30\,\% de jumeaux
monozygotes (appelés \og  vrais jumeaux \fg{} qui sont obligatoirement de même sexe : deux garçons ou deux filles) et donc 70\,\% de jumeaux dizygotes (appelés \og  faux jumeaux \fg, qui
peuvent être de sexes différents: deux garçons, deux filles ou un garçon et une fille).

Dans le cas de naissances doubles, on admet que, comme pour les naissances ordinaires, la probabilité d'être une fille à la naissance est égale à $0,49$ et que celle d'être un garçon à la
naissance est égale à $0,51$.

Dans le cas d'une naissance double de jumeaux dizygotes, on admet aussi que le sexe du
second nouveau--né des jumeaux est indépendant du sexe du premier nouveau-né.

On choisit au hasard un accouchement double réalisé dans cette maternité et on considère les évènements suivants :

\begin{itemize}
\item[$\bullet~~$]$M$ : \og  les jumeaux sont monozygotes \fg{} ;
\item[$\bullet~~$]$F_1$ : \og  le premier nouveau--né est une fille \fg{} ;
\item[$\bullet~~$]$F_2$ : \og  le second nouveau--né est une fille \fg.
\end{itemize}

\emph{On notera $p(A)$ la probabilité de l'évènement $A$ et $\overline{A}$ l'évènement contraire de $A$.}

\begin{center}
\pstree[treemode=R,nodesepA=0pt,nodesepB=2.5pt,treesep = 1cm,levelsep=2.5cm]{\TR{}}
{\pstree{\TR{$M$~~}\naput{0,3}}
	{\pstree{\TR{$F_1$}\naput{0,49}}
		{\TR{$F_2$}\naput{1}	
		}
	\pstree{\TR{$\overline{F_1}$~}\naput{0,51}}
		{\TR{$\overline{F_2}$}\nbput{1}	
		}
	}
\pstree{\TR{$\overline{M}$~~}\nbput{0,7}}
	{\pstree{\TR{$F_{1}$~}\naput{0,49}}
		{\TR{$F_2$}\naput{0,49}
		\TR{$\overline{F_2}$}\nbput{0,51}
		}
	\pstree{\TR{$\overline{F_1}$~}\naput{0,51}}
		{\TR{$F_2$}\naput{0,49}
		\TR{$\overline{F_2}$}\nbput{0,51}
		}
	}
}
\end{center}
\smallskip

	\begin{enumerate}
		\item 
Voir ci-dessus
		\item 
On a en suivant les première et troisième branche de l'arbre pondéré :

$p\left(F_1 \cap F_2 \right) = 0,3 \times 0,49 \times 1 + 0,7 \times 0,49 \times 0,49 = 0,147 + \np{0,16807} = \np{0,31507} \approx 0,316$.
		\item 
Sachant que les nouveaux-nés sont des jumelles, la probabilité qu'elles soient monozygotes est la probabilité conditionnelle :

$p_{F_1 \cap F_2} (M) = \dfrac{p(M \cap (F_1 \cap F_2))}{p(F_1 \cap F_2)} = \dfrac{0,3 \times 0,49 \times 1}{\np{0,31507}} \approx \np{0,4665}$ soit 0,467 au millième près.
	\end{enumerate}
\end{enumerate}

\bigskip

\end{document}
